% !PRINTED 446
\weberpage{446}
But by \S.\,115, (14) and (8) we have

\[
\alpha^2 + \beta^2\varepsilon^3 + \alpha\,\beta\,\varepsilon^2 = \alpha^2 + \beta\,(\alpha\,\varepsilon^2 + \beta\,\varepsilon^3) = \alpha^2 - \beta\,\gamma = -\,\alpha,
\]
and if one transforms the three coefficients of the expression (8) in this way, there results

\[
-\,\alpha\,\gamma\,x_1 - \beta\,\gamma\,\varepsilon^3 x_2 - \alpha\,\beta\,\varepsilon^2 x_3,
\]
that is, $A$ changes its sign under application of the substitution $\omega$, and the principal axis $A$ therefore remains unchanged by the substitution $\omega$.

Since $A$ moreover remains unchanged by the substitution $\Theta$, because two poles of $\Theta$ lie on $A$, we have the whole group of the 8th degree by which $A$ remains unchanged, which we call the \gesperrt{group of} $A$ and denote by $G_a$, in the form

\weq{9}{\omega^u\,\Theta^v.}
Among these eight substitutions there are only two, namely the identical one and $\Theta^2$, which leave all points of the line $A$ unchanged. Under the others only two points each remain at rest. For if one calculates, by the path followed here, from \S.\,115, (5), (16) the principal axes of $\omega$, $\omega\Theta$, $\omega\Theta^2$, $\omega\Theta^3$, one obtains, in the notation (7),

\[
A_{2,1},\ A_{3,0},\ A_{3,3},\ A_{2,0},
\]
which are all different from the principal axis $A_{1,0}$ of $\Theta^2$.

Besides the principal axis $A$, the substitution $\Theta^2$ has yet an isolated pole, which we obtain when we choose the root $\lambda = 1$ of the cubic equation (4). Then for the coordinates of this pole there result the equations

\[
\begin{array}{l}
(\beta-1)\,x_1 + \gamma\,\varepsilon^3 x_2 + \alpha\,\varepsilon^2 x_3 = 0,\\
\gamma\,\varepsilon^{-3} x_2 + (\alpha-1)\,x_2 + \beta\,\varepsilon^{-1} x_3 = 0,\\ % [sic] printed x_2 twice, expected x_1 in first term
\alpha\,\varepsilon^{-2} x_1 + \beta\,\varepsilon\, x_2 + (\gamma-1)\,x_3 = 0,
\end{array}
\]
and from these one obtains

\weq{10}{x_1 : x_2 : x_3 = \alpha\,\gamma\,\varepsilon^4 : \beta\,\gamma\,\varepsilon : \alpha\,\beta\,\varepsilon^2.}
Of such points there are 21. The point (10) we denote for the moment by $a$. Now the point $a$ remains unchanged not only by $\Theta$, but also, as a calculation on the basis of the formulae \S.\,115, (17), (14) shows, by $\omega$, and consequently by the whole group $G_a$, and is therefore a pole of order at least eight. But since it can be brought by 21 substitutions into 21 different positions, it is also not \gesperrt{more} than
%CONT>>
